problem:  
Let \( (M^n, g) \) be a compact, simply connected Riemannian manifold with **quasipositive curvature**, meaning \( \sec \ge 0 \) everywhere and \( \sec > 0 \) at least at one point. Suppose \( M \) admits an **effective** isometric torus action \( T^r \curvearrowright M \) of rank \( r = \lfloor \tfrac{n+1}{2} \rfloor \), the maximal possible rank for a positively curved manifold.  

**Question:**  
Must \( M \) be diffeomorphic to a compact rank-one symmetric space (CROSS)?  
Equivalently, does maximal torus symmetry rank together with quasipositive curvature imply that \( M \) is diffeomorphic to one of the known examples \( S^n, \mathbb{C}P^{n/2}, \mathbb{H}P^{n/4}, \mathrm{CaP}^2 \)?  

Why is it a "good" problem:  
This problem explicitly extends a cornerstone result in Riemannian geometry—the **Grove–Searle symmetry rank theorem**, which asserts that if \( M^n \) has **positive** sectional curvature and an isometric torus action of maximal rank \( \lfloor (n+1)/2 \rfloor \), then \( M \) must be a CROSS. The present question asks whether this rigidity persists when the curvature is only **quasipositive**, that is, positive somewhere but allowed to vanish elsewhere.  

This is a “good” problem because:  

- **Profound insight:** It investigates whether the phenomenon of **maximal symmetry enforcing global geometric and topological rigidity** survives when the curvature assumption is relaxed. This probes the boundary between the known rigid world of positive curvature and the largely mysterious world of nonnegative or quasipositive curvature.  

- **Bridges multiple fields:** Resolving this question would require combining techniques from **transformation group theory**, **Alexandrov geometry**, and **metric collapse theory**, possibly using methods from **Ricci flow** or **topological symmetry analysis**. The problem sits precisely at the interface of curvature pinching and orbit structure classification.  

- **Conceptual simplicity but deep difficulty:** The statement is strikingly simple: “Does maximal symmetry force a CROSS under quasipositive curvature?” Yet its resolution requires understanding how local positive curvature behavior may or may not propagate globally under strong symmetry constraints—a major unsolved issue.  

- **New geometric challenges:** In the quasipositive case, flat regions and degenerate orbits can obstruct standard comparison theorems and convexity arguments crucial in the Grove–Searle proof. Determining whether one point of strict curvature positivity is sufficient to enforce global rigidity remains beyond current techniques.  

Thus, this question naturally extends a classical rigidity theorem into an unexplored curvature regime, linking local geometry, global topology, and transformation group structure—making it a genuinely **good research problem**.